%% ch-recipes.tex — FORTH Cookbook Recipes for daily use
%% SOURCE: .claude/CLAUDE.md, src/word_source/
%% Included from user-guide/main.tex

%%━━━━━━━━━━━━━━━━━━━━━━━━━━━━━━━━━━━━━━━━━━━━━━━━━━━━━━━━━━━━━━━━━━━━━━━━━
\part{Recipes}
%%━━━━━━━━━━━━━━━━━━━━━━━━━━━━━━━━━━━━━━━━━━━━━━━━━━━━━━━━━━━━━━━━━━━━━━━━━

%%────────────────────────────────────────────────────────────────────────────
\chapter{Frequently Used Idioms}
\label{chap:idioms}
%%────────────────────────────────────────────────────────────────────────────

This chapter collects short FORTH idioms you will reach for repeatedly. Each
entry shows the idiom, its stack effect, and a concrete example. They are
organized by what you are trying to accomplish, not by word name.

%% ─────────────────────────────────────────────────────────────────────────
\section{Swap and Drop a Pair}
\label{sec:idiom-nip}

Discard the second item from the top:

\begin{lstlisting}[language=Forth]
( a b -- b )
SWAP DROP         \ explicit
NIP               \ compact equivalent
\end{lstlisting}

Use \texttt{NIP} when you compute something from two values but only need the
result of the second operation:

\begin{lstlisting}[language=Forth]
: AREA-FROM-SIDES  ( width height -- area )
    * ;   \ stack: width height -- area; no NIP needed here
\end{lstlisting}

\begin{lstlisting}[language=Forth]
: LAST-OF-THREE  ( a b c -- c )
    NIP NIP ;
\end{lstlisting}

%% ─────────────────────────────────────────────────────────────────────────
\section{Apply an Operation to Two Copies}
\label{sec:idiom-tuck}

\texttt{TUCK} lets you use a value in two consecutive operations without
re-fetching it from a variable:

\begin{lstlisting}[language=Forth]
( n -- n*n  n )      \ square, then original
TUCK *
\end{lstlisting}

\begin{lstlisting}[language=Forth]
( a b -- b a b )
TUCK
\end{lstlisting}

Use \texttt{TUCK} when you need the top value both as an operand and to keep it
on the stack:

\begin{lstlisting}[language=Forth]
: PERCENT-OF  ( total pct -- value remaining )
    \ pct of total, then total minus that value
    TUCK  100 */ DUP ROT SWAP - ;
\end{lstlisting}

%% ─────────────────────────────────────────────────────────────────────────
\section{Conditional Increment}
\label{sec:idiom-cond-inc}

Increment a counter only if a condition is true. FORTH's flag convention
(\texttt{-1} for true, \texttt{0} for false) means you can add a flag
directly:

\begin{lstlisting}[language=Forth]
( counter condition -- counter+1|counter )
+        \ adds -1 (true) or 0 (false) -- this subtracts 1!
SWAP -   \ correct: condition is -1, negating gives 1
\end{lstlisting}

More clearly:

\begin{lstlisting}[language=Forth]
: COND-INC  ( n flag -- n' )
    IF  1 +  THEN ;
\end{lstlisting}

\begin{lstlisting}[language=Forth]
ok> 5 1 3 > COND-INC .    \ 1 > 3 is false (0), so no increment
5 ok>
ok> 5 3 1 > COND-INC .    \ 3 > 1 is true (-1), so increment
6 ok>
\end{lstlisting}

%% ─────────────────────────────────────────────────────────────────────────
\section{Clamp a Value to a Range}
\label{sec:idiom-clamp}

\begin{lstlisting}[language=Forth]
: CLAMP  ( n lo hi -- n' )
    ROT     \ ( lo hi n )
    MIN     \ ( lo min(hi,n) )
    MAX ;   \ ( max(lo, min(hi,n)) )
\end{lstlisting}

\begin{lstlisting}[language=Forth]
ok> 15 0 10 CLAMP .
10 ok>
ok> -5 0 10 CLAMP .
0 ok>
ok>  7 0 10 CLAMP .
7 ok>
\end{lstlisting}

%% ─────────────────────────────────────────────────────────────────────────
\section{Absolute Difference}
\label{sec:idiom-abs-diff}

\begin{lstlisting}[language=Forth]
: ABS-DIFF  ( a b -- |a-b| )
    - ABS ;
\end{lstlisting}

%% ─────────────────────────────────────────────────────────────────────────
\section{Integer Division with Rounding}
\label{sec:idiom-round-div}

FORTH's \texttt{/} truncates toward zero. For rounded division:

\begin{lstlisting}[language=Forth]
: ROUND-DIV  ( a b -- a/b rounded )
    2DUP /        \ quotient
    -ROT MOD      \ remainder  ( quot rem )
    SWAP OVER     \ ( rem quot rem )
    2/ ABS > IF   \ remainder > |divisor/2|?
        1 +       \ round up
    THEN ;
\end{lstlisting}

Or, more simply for positive values:

\begin{lstlisting}[language=Forth]
: ROUND-DIV+  ( a b -- rounded )
    2DUP  2/  +  SWAP / ;   \ (a + b/2) / b
\end{lstlisting}

%% ─────────────────────────────────────────────────────────────────────────
\section{Swap Top Two Booleans}
\label{sec:idiom-bool-logic}

Because FORTH booleans are plain integers, boolean operations are just
arithmetic. \texttt{NOT} inverts, \texttt{AND} is bitwise AND (works correctly
on \texttt{-1}/\texttt{0} flags), and \texttt{OR} is bitwise OR:

\begin{lstlisting}[language=Forth]
ok> -1 0 AND .     \ true AND false = false
0 ok>
ok> -1 0 OR  .     \ true OR false = true
-1 ok>
ok> -1 NOT   .     \ NOT true = false
0 ok>
\end{lstlisting}

%% ─────────────────────────────────────────────────────────────────────────
\section{String Equality Test}
\label{sec:idiom-str-eq}

Compare two strings of known length:

\begin{lstlisting}[language=Forth]
: STR=  ( addr1 len1 addr2 len2 -- flag )
    ROT OVER <> IF  2DROP DROP  0  EXIT  THEN
    \ lengths are equal; now compare bytes
    OVER SWAP  ( addr1 addr2 len )
    0 DO
        OVER C@  OVER C@  <> IF
            2DROP  0  UNLOOP  EXIT
        THEN
        SWAP 1+ SWAP 1+
    LOOP
    2DROP  -1 ;
\end{lstlisting}

%% ─────────────────────────────────────────────────────────────────────────
\section{Fibonacci Sequence}
\label{sec:idiom-fib}

Classic Fibonacci using the return stack as a third register:

\begin{lstlisting}[language=Forth]
: FIB  ( n -- fib(n) )
    DUP 2 < IF EXIT THEN   \ base cases: fib(0)=0, fib(1)=1
    DUP 1 - RECURSE        \ fib(n-1)
    SWAP 2 - RECURSE       \ fib(n-2)
    + ;
\end{lstlisting}

\begin{lstlisting}[language=Forth]
ok> 10 FIB .
55 ok>
\end{lstlisting}

For larger values, use the iterative form (much faster):

\begin{lstlisting}[language=Forth]
: FIB-ITER  ( n -- fib(n) )
    DUP 2 < IF EXIT THEN
    1 0           \ a=1, b=0 (fib(1), fib(0))
    ROT 1 - 0 DO
        OVER +    \ new a = a+b
        SWAP      \ rotate
    LOOP
    DROP ;
\end{lstlisting}

\begin{lstlisting}[language=Forth]
ok> 40 FIB-ITER .
102334155 ok>
\end{lstlisting}

%% ─────────────────────────────────────────────────────────────────────────
\section{Printing a Separator Line}
\label{sec:idiom-separator}

\begin{lstlisting}[language=Forth]
: HLINE  ( n -- )
    0 DO  45 EMIT  LOOP  CR ;   \ 45 = '-'

: BANNER  ( addr len -- )
    DUP 2 + >R
    ." +-"  R@ HLINE  ." | "
    2DUP TYPE  CR  ." +-"  R> HLINE ;
\end{lstlisting}

\begin{lstlisting}[language=Forth]
ok> S" StarForth v3.1" BANNER
+----------------+
| StarForth v3.1
+----------------+
ok>
\end{lstlisting}

%%────────────────────────────────────────────────────────────────────────────
\chapter{Working Examples}
\label{chap:examples}
%%────────────────────────────────────────────────────────────────────────────

%% ─────────────────────────────────────────────────────────────────────────
\section{A Simple Stack Calculator}
\label{sec:example-calculator}

This example builds an interactive RPN calculator on top of the StarForth
REPL. Because StarForth IS an RPN calculator, this is mostly a demonstration
of how you can add user-level commands on top:

\begin{lstlisting}[language=Forth]
\ A FORTH-based RPN calculator session

\ Define helper display words
: SHOW-STACK  .S CR ;
: RESULT      . CR ;

\ Factorial
: FACT  ( n -- n! )
    DUP 0= IF  DROP 1  EXIT  THEN
    DUP 1 - RECURSE  * ;

\ Power
: POW  ( base exp -- base^exp )
    1 SWAP 0 DO OVER * LOOP  NIP ;
\end{lstlisting}

Using these at the REPL:

\begin{lstlisting}[language=Forth]
ok> 5 FACT RESULT
120
ok> 2 10 POW RESULT
1024
ok> 12 FACT RESULT
479001600
\end{lstlisting}

%% ─────────────────────────────────────────────────────────────────────────
\section{A Data Table}
\label{sec:example-data-table}

Store and retrieve structured data using \texttt{CREATE} and array indexing:

\begin{lstlisting}[language=Forth]
\ A table of (name-hash, value) pairs
\ 5 entries, 2 cells each

CREATE DATA-TABLE 10 CELLS ALLOT

: TABLE-SET  ( value index -- )
    2* CELLS DATA-TABLE + ! ;

: TABLE-GET  ( index -- value )
    2* CELLS DATA-TABLE + @ ;

\ Initialize
100  0 TABLE-SET
200  1 TABLE-SET
300  2 TABLE-SET
400  3 TABLE-SET
500  4 TABLE-SET

: DUMP-TABLE
    ." Index  Value" CR
    5 0 DO
        I  4 .R  I TABLE-GET  8 .R  CR
    LOOP ;
\end{lstlisting}

\begin{lstlisting}[language=Forth]
ok> DUMP-TABLE
Index  Value
    0    100
    1    200
    2    300
    3    400
    4    500
ok>
\end{lstlisting}

%% ─────────────────────────────────────────────────────────────────────────
\section{A Simple Moving Average}
\label{sec:example-moving-average}

Compute a moving average over the last N readings using a circular buffer:

\begin{lstlisting}[language=Forth]
10 CONSTANT MA-SIZE
CREATE MA-BUF   MA-SIZE CELLS ALLOT
VARIABLE MA-INDEX
VARIABLE MA-SUM

0 MA-INDEX !
0 MA-SUM !
MA-BUF  MA-SIZE CELLS  0 FILL   \ zero the buffer

: MA-UPDATE  ( new-value -- average )
    \ subtract old value at current index
    MA-INDEX @ CELLS MA-BUF + @
    MA-SUM @ SWAP - MA-SUM !
    \ store new value
    DUP MA-INDEX @ CELLS MA-BUF + !
    \ add to sum
    MA-SUM @ + MA-SUM !
    \ advance index circularly
    MA-INDEX @ 1 + MA-SIZE MOD MA-INDEX !
    \ return average
    MA-SUM @ MA-SIZE / ;

: TEST-MA
    ." Moving average test (window=10)" CR
    20 0 DO
        I 3 *    \ input: 0,3,6,9,...,57
        MA-UPDATE
        ." Input: " I 3 * 4 .R
        ."  Avg: " . CR
    LOOP ;
\end{lstlisting}

%% ─────────────────────────────────────────────────────────────────────────
\section{A Minimal State Logger}
\label{sec:example-logger}

Log timestamped events to a block buffer:

\begin{lstlisting}[language=Forth]
\ Minimal event logger using block 2048
2048 CONSTANT LOG-BLOCK
0 CONSTANT LOG-OFFSET   \ offset within block (in bytes)
VARIABLE LOG-POS
0 LOG-POS !

: LOG-RESET
    0 LOG-POS !
    LOG-BLOCK BUFFER  1024 0 FILL  UPDATE ;

: LOG-WRITE  ( addr count -- )
    LOG-POS @  1024 <  IF
        LOG-BLOCK BUFFER LOG-POS @ +
        SWAP MOVE
        LOG-POS @ + LOG-POS !
    ELSE
        ." Log full" CR
    THEN ;

: LOG-MSG  ( addr count -- )
    2DUP LOG-WRITE  DROP DROP ;

: LOG-DUMP
    LOG-BLOCK BLOCK  LOG-POS @ TYPE CR ;
\end{lstlisting}

%% ─────────────────────────────────────────────────────────────────────────
\section{Physics-Aware Benchmarking}
\label{sec:example-physics-bench}

A benchmark harness that controls the adaptive runtime for reproducible
measurements:

\begin{lstlisting}[language=Forth]
\ Benchmark framework using physics controls

VARIABLE BENCH-ITERATIONS
100000 BENCH-ITERATIONS !

: BENCH-WORD  ( xt -- )
    \ Warm up (let physics settle)
    PHYSICS-RESET-STATS
    1000 0 DO  DUP EXECUTE  LOOP
    \ Now freeze and measure
    PHYSICS-FREEZE
    BENCH-ITERATIONS @ 0 DO  DUP EXECUTE  LOOP
    PHYSICS-THAW
    DROP ;

\ Example: benchmark DUP
: RUN-DUP-BENCH
    ." Benchmarking DUP x "
    BENCH-ITERATIONS @ . ." iterations" CR
    ' DUP  BENCH-WORD
    ." Done." CR ;
\end{lstlisting}

The \texttt{PHYSICS-FREEZE}/\texttt{PHYSICS-THAW} pair ensures that heat
updates, cache evictions, and decay events do not occur during the timed
region, giving cleaner timing measurements.

%%────────────────────────────────────────────────────────────────────────────
\chapter{Quick Reference: Key Sequences}
\label{chap:key-sequences}
%%────────────────────────────────────────────────────────────────────────────

%% TODO(bob): verify which readline / terminal sequences StarForth actually
%% supports in the REPL — these are standard terminal behaviors, not FORTH words.

The StarForth REPL runs in a standard POSIX terminal. The following key
sequences are handled by the terminal or readline layer:

\begin{center}
\begin{tabular}{ll}
\toprule
Key sequence & Effect \\
\midrule
\textbf{Enter} & Execute the current input line \\
\textbf{Ctrl-C} & Interrupt current operation; return to prompt \\
\textbf{Ctrl-D} & End of input (exits REPL, same as \texttt{BYE}) \\
\textbf{Backspace} & Delete previous character \\
\textbf{Up arrow} & Previous history entry \\
\textbf{Down arrow} & Next history entry \\
\textbf{Ctrl-A} & Move cursor to beginning of line \\
\textbf{Ctrl-E} & Move cursor to end of line \\
\textbf{Ctrl-K} & Kill to end of line \\
\textbf{Ctrl-U} & Kill entire line \\
\bottomrule
\end{tabular}
\end{center}

%% TODO(bob): confirm whether StarForth uses GNU readline or its own line
%% editor. Add a section on the block editor key bindings once confirmed.

%%────────────────────────────────────────────────────────────────────────────
\chapter{Migrating from Other FORTH Systems}
\label{chap:migrating}
%%────────────────────────────────────────────────────────────────────────────

%% ─────────────────────────────────────────────────────────────────────────
\section{From gforth or ANS FORTH}
\label{sec:migrate-gforth}

StarForth targets FORTH-79 compliance. It is broadly compatible with ANS FORTH
(FORTH-94) and gforth programs, with a few differences to watch:

\begin{description}
  \item[\texttt{S"} handling] StarForth's \texttt{S"} behaves as in ANS FORTH:
    it pushes an address and count. Programs that assume counted-string format
    (a length byte followed by characters) need to use \texttt{COUNT} to
    convert.

  \item[Character size] ANS FORTH specifies \texttt{CHAR} as one address unit.
    StarForth uses byte-addressed cells, so \texttt{CHAR+} is \texttt{1+} and
    \texttt{CHARS} is a no-op. Code that uses \texttt{CHARS} for portability
    will work correctly.

  \item[Cell size] StarForth cells are 64 bits on a 64-bit host. Programs
    that assume 32-bit cells will need adjustment if they use bit patterns that
    fill a cell exactly.

  \item[LOCAL variables] StarForth does not implement ANS FORTH local variables
    (\texttt{LOCALS|} / \texttt{TO}). Use named \texttt{VARIABLE}s instead, or
    restructure to keep values on the stack.

  \item[THROW/CATCH] StarForth's exception handling may differ from ANS FORTH.
    %% TODO(bob): confirm whether THROW/CATCH are implemented and to what spec.

  \item[Floating-point] StarForth does not implement the ANS FORTH floating-point
    word set. Use Q48.16 fixed-point (\texttt{Q+}, \texttt{Q*}, etc.) or integer
    arithmetic for fractional values.
\end{description}

%% ─────────────────────────────────────────────────────────────────────────
\section{From FORTH-79}
\label{sec:migrate-forth79}

StarForth is FORTH-79 compliant. Programs written for the original FORTH-79
standard should run without modification. The main additions StarForth makes
beyond FORTH-79 are:

\begin{itemize}
  \item The physics adaptive runtime (seven feedback loops), which is entirely
    transparent unless you call \texttt{WORD-ENTROPY},
    \texttt{PHYSICS-FREEZE}, or \texttt{PHYSICS-THAW}.
  \item The ACL word-level access control system, disabled by default.
  \item Q48.16 fixed-point arithmetic words.
  \item The block subsystem with pluggable backends.
  \item StarForth-specific extensions in \texttt{starforth\_words.c}.
\end{itemize}

None of these conflict with FORTH-79 code. They simply add capability on top.

%% ─────────────────────────────────────────────────────────────────────────
\section{From C or Python}
\label{sec:migrate-c-python}

The largest adjustment is postfix notation and the stack discipline. Some
practical advice:

\begin{itemize}
  \item Draw a picture of the stack before each operation while learning.
    Writing \texttt{( a b -- c )} for every sub-expression helps.

  \item Use \texttt{.S} compulsively at first. It does not consume anything,
    so you can sprinkle it everywhere.

  \item Resist the urge to write long, complex word bodies. Five lines is
    often too many. Factor into smaller words.

  \item \texttt{VARIABLE} fills the role of a local variable but it is global
    within the VM --- be careful with re-entrancy if a word using a
    \texttt{VARIABLE} calls itself or a word that also uses the same variable.

  \item There are no types. You decide what each stack cell means. A cell could
    be an integer, an address, a boolean, a character, or a \Qtype{} value.
    Naming your stack-effect slots meaningfully (\texttt{( count addr -- result )})
    is your type system.
\end{itemize}
